\section{Metric Spaces}

\begin{Definition}{Metric Space}{metric-space}
    Let $X$ be a non-empty set. A \textbf{metric}\footnote{\href{https://trello.com/c/3EPccNTa}{Introduction to Topology and Functional Analysis by Goerge F. Simmons}, p.51}
    on $X$ is a real function $d$ of ordered pairs of elements of $X$ which satisfies the following conditions

    \begin{itemize}
        \item $d(x, y) \ge 0$, and $d(x, y) = 0 \iff x = y$
        \item $d(x, y) = d(y, x)$ (symmetry)
        \item $d(x, y) \le d(x, z) + d(z, y)$ (the triangle inequality)
    \end{itemize}
\end{Definition}

$d(x, y)$ is called the \textit{distance} between $x$ and $y$. Thus a metric space consists of 2 objects:

\begin{enumerate}
    \item a non-empty set $X$, and
    \item a metric $d$ on $X$
\end{enumerate}

\footnote{\href{https://trello.com/c/3EPccNTa}{Introduction to Topology and Functional Analysis by Goerge F. Simmons}, p.70}
Now let $X$ be a metric space with metric $d$, and let

\begin{equation}
    \{ x_n \} = \{ x_1, x_2, \cdots, x_n, \cdot \}
\end{equation}

be a sequence of points in $X$. We say that $\{ x_n \}$ is \textbf{convergent} if there exists a point $x$ in $X$ such
that either

\begin{enumerate}
    \item for each $\epsilon > 0$, there exists a positive integer $n_0$ such that $n \ge n_0 \Rightarrow d(x_n, x) < \epsilon$,
          or, equivalently,
    \item for each open sphere $S_{\epsilon}(x)$ centered on $x$, there exists a positive integer $n_0$ such that $x_n$
          is in $S_{\epsilon}(x)$ for all $n \ge n_0$
\end{enumerate}

We say in this case that $x_n$ converges to $x$ and $x$ is called the \textit{limit} of the sequence $\{ x_n \}$ and we
sometimes write $x_n \rightarrow x$ in the form of

\[ \lim x_n = x \]

Now let's consider two cases of convergence:

\begin{itemize}
    \item for each $\epsilon > 0$, there exists a positive integer $n_0$ such that $n \ge n_0 \Rightarrow d(x_n, x) < \epsilon$ (the same definition above)
    \item for each $\epsilon > 0$, there exists a positive integer $n_0$ such that $m \ge n_0 \Rightarrow d(x_m, x) < \epsilon$
\end{itemize}

By the \textit{symmetry} and \textit{triangle inequality} of the Def.~\ref{def:metric-space}:

\begin{equation}
    d(x_m, x_n) \le d(x_m, x) + d(x, x_n) = d(x_m, x) + d(x_n, x) \le \frac{\epsilon}{2} + \frac{\epsilon}{2} = \epsilon
\end{equation}

for all $m, n \lg n_0$. Therefore, every convergent sequence $\{ x_n \}$ has the following property:

\begin{tcolorbox}[
    colback=green!5!white,
    colframe=green!75!black,
    colbacktitle=red!85!black
]
    \begin{center}
        For each $\epsilon > 0$, there exists a positive integer $n_0$ such that $m, n \ge n_0 \Rightarrow d(x_m, x_n) < \epsilon$
    \end{center}
\end{tcolorbox}

A sequence with this property is called a \roundinlinebox[green]{Cauchy sequence}. Intuitively, a Cauchy sequence is a
sequence whose elements become arbitrarily close to each other as the sequence progress:

\begin{center}
    \begin{tikzpicture}
        \draw[->] (-0.2,0) -- (4.5,0) node[right] {$n$};
        \draw[->] (0,-1.2) -- (0,2.5) node[above] {$x_n$};
        \draw[dash pattern=on 0pt off 2\pgflinewidth,line cap=round,line width=0.5mm,color=blue,domain=0:4,samples=100]  plot (\x,{1+ sin(10*\x r)*e^(-\x)});
        \draw[dash pattern=on 0pt off 2\pgflinewidth,line cap=round,line width=0.5mm,color=red,dashed,thin,domain=0:4,samples=100]  plot (\x,{1+e^(-\x)});
        \draw[dash pattern=on 0pt off 2\pgflinewidth,line cap=round,line width=0.5mm,color=red,dashed,thin,domain=0:4,samples=100]  plot (\x,{1-e^(-\x)});
    \end{tikzpicture}
\end{center}

In addition, we have also shown that every convergent sequence is a Cauchy sequence. \textit{The converse of this,
however, is not necessarily true}. That is, a Cauchy sequence is not necessarily convergent. As an example, consider the
subspace $X = (0, 1]$ of the real line where $x_n \in X$. The sequence defined by $x_n = \frac{1}{n}$ is easily seen to
be a Cauchy sequence in this space, but it is not convergent, because the point $0 \not\in X$ is not a point of the
space

\begin{figure}[H]
    \centering
    \begin{tikzpicture}
        \colorlet{mygreen}{green!70!black}

        \draw[->] (-0.2,0) -- (10,0) node[right] {$n$};
        \draw[->] (0,-1.2) -- (0,3) node[above] {$x_n$};
        \draw[dash pattern=on 0pt off 2\pgflinewidth,line cap=round,line width=0.5mm,color=blue,domain=0.3:1,samples=20]  plot (\x,{\x^(-1)});
        \draw[line width=0.5mm,color=blue,domain=1:10]  plot (\x,{\x^(-1)});

        \coordinate (A) at (0, 1);
        \coordinate (B) at (0, 0);
        \fill[mygreen] (A) circle (1.5pt) node[left=1] {$x_n = 1 \in X$};
        \fill[red] (B) circle (1.5pt) node[left=1] {$x_n = 0 \not\in X$};

        \draw [dashed] (A) -- (1, 1);
    \end{tikzpicture}
    \caption{Here, the real numbers in the range of $X = (0, 1]$ takes on the values ``y-axis". We use a solid line to
    denote the range of $n$ values under the subspace of $X$. It really should be a dotted line because $n$ can only be
    integers. We do that in order to distinguish it from the range we are not concerned with, i.e. $X = (1, \infty)$}
\end{figure}

\highlight[green]{The notion of a convergent sequence is not intrinsic to the sequence itself, but also depends on
the structure of the space in which it lies}. A convergent sequence is not convergent "on its own"; it must converge to
some point in the space.

Sounds wordy? When the space \href{https://www.reddit.com/r/learnmath/comments/exp735/comment/fgaqd8p/?utm\_source=share&utm\_medium=web3x&utm\_name=web3xcss&utm\_term=1&utm\_content=share\_button}{has a hole} in it. Like if you remove 0 from $\mathbb{R}$ like in the case above, then
$\frac{1}{n}$ is a Cauchy sequence that doesn't converge.

We are now ready to introduce the concept of \textit{complete metric space}, which we define as follows:

\begin{Definition}{Complete Metric Space}{complete-metric-space}
    A complete metric space is a metric space in which every Cauchy sequence is convergent
\end{Definition}

\begin{Definition}{Linear Space}{linear-space}
    Let $L$ be a non-empty set, and assume that each pair of elements $x$ and $y$ in $L$ can be combined by a process
    called \textbf{addition} to yield an element $z$ in $L$ denoted by $z = x + y$. Assume also that this operation of
    addition satisfies the following conditions:

    \begin{itemize}
        \item $x + y = y + x$
        \item $x + (y + z) = (x + y) + z$
        \item There exists in $L$ a unique element, denoted by $0$ and called the \textbf{zero element}, or the origin,
              such that $x + 0 = x$ for every $x$
        \item To each element $x$ in $L$ there corresponds a unique element in $L$, denoted by $-x$ and called the
              negative of $x$, such that $x + (-x) = 0$.
    \end{itemize}

    We adopt the device of referring to the system of real numbers or to the system of complex numbers as the
    \textbf{scalers}. We now assume that each scalar $\alpha$ and each element $x$ in $L$ can be combined by a process
    called \textbf{scalar multiplication} to yield an element $y$ in $L$ denoted by $y = \alpha x$ in such a way that

    \begin{itemize}
        \item $\alpha(x + y) = \alpha x + \alpha y$
        \item $(\alpha + \beta)x = \alpha x + \beta x$
        \item $(\alpha\beta)x = \alpha(\beta x)$
        \item $1 \cdot x = x$
    \end{itemize}

    The albegraic system $L$ defined by these operations and axioms is called a
    \textbf{linear space}\footnote{\href{https://trello.com/c/3EPccNTa}{Introduction to Topology and Functional Analysis by Goerge F. Simmons}, p.81}
\end{Definition}

Depending on the numbers admitted as scalars (only the real numbers, or all the complex numbers), we distinguish when
necessary between \textit{real linear spaces} and \textit{complex linear spaces}. A linear space is often called a
\hyperlink{vector-space}{\textit{vector space}} and its elements are spoken of as \textit{vectors}.

\begin{Definition}{Normed Linear Space}{normed-linear-space}
    A \textbf{Normed Linear Space}\footnote{\href{https://trello.com/c/3EPccNTa}{Introduction to Topology and Functional Analysis by Goerge F. Simmons}, p.81}
    is a linear space on which there is a \textbf{norm} defined, i.e. a function which assigns to each element $x$ in
    the space a real number $\Vert x \Vert$ in such a manner that

    \begin{itemize}
        \item $\Vert x \Vert \ge 0$, and $\Vert x \Vert = 0 \iff x = 0$
        \item $\Vert x + y \Vert \le \Vert x \Vert + \Vert y \Vert$
        \item $\Vert \alpha x \Vert = \vert \alpha \vert \Vert x \Vert$
    \end{itemize}
\end{Definition}

Intuitively, a normed linear space is simply a linear space in which a notion of the distance from an arbitrary element
to origin is defined.

\begin{Definition}{Banach space}{banach-space}
    A \textbf{Banach space}\footnote{\href{https://trello.com/c/3EPccNTa}{Introduction to Topology and Functional Analysis by Goerge F. Simmons}, p.81}
    is a normed linear space which is complete as a metric space.
\end{Definition}



One of the principal applications of the tory of Banach algebras develped from Def.~\ref{def:metric-space} to
Def~\ref{def:banach-space} is to the study of operators in Hilbert
spaces.\footnote{\href{https://trello.com/c/3EPccNTa}{Introduction to Topology and Functional Analysis by Goerge F. Simmons}, p.243}

\begin{Definition}{Hilbert Space}{hilbert-space}
    A \textbf{Hilbert space}\footnote{\href{https://trello.com/c/3EPccNTa}{Introduction to Topology and Functional Analysis by Goerge F. Simmons}, p.245}
    is a complex Banach space whose norm arises from an inner product in which there is defined a complex function
    $(\boldsymbol{v}, \boldsymbol{w})$ of vectors $\boldsymbol{v}$ with the following properties:

    \begin{itemize}
        \item $(\alpha \boldsymbol{v} + \beta \boldsymbol{w}, \boldsymbol{u}) = \alpha(\boldsymbol{v}, \boldsymbol{u}) + \beta(\boldsymbol{w}, \boldsymbol{u})$
        \item $\overline{(\boldsymbol{v}, \boldsymbol{w})} = (\boldsymbol{w}, \boldsymbol{v})$
        \item $(\boldsymbol{v}, \boldsymbol{v}) = \Vert x \Vert^2$
    \end{itemize}

    \tcblower

    A complex function\footnote{\href{https://trello.com/c/3EPccNTa}{Introduction to Topology and Functional Analysis by Goerge F. Simmons}, p.17}
    is one whose range consists of complex numbers

    If $z = a + ib$ is a complex number, then its conjugate $\overline{z}$ is defined by $\overline{z} = a + i(-b)$\footnote{\href{https://trello.com/c/3EPccNTa}{Introduction to Topology and Functional Analysis by Goerge F. Simmons}, p.53}
\end{Definition}

\section{Vector Spaces and Subspaces}

\begin{Definition}{Vector Space \phantomsection\hypertarget{vector-space}}{vector-space}
    The space $\mathbb{R}^n$ consists of all column vectors $\boldsymbol{v}$ with $n$ components\footnote{\href{https://trello.com/c/qHJeDNkU}{Introduction to Linear Algebra, Strang, 4th Edition, 2009}, p. 120}
\end{Definition}

The components of $\boldsymbol{v}$ are real numbers, which is the reason for the letter $\mathbb{R}$. A vector whose $n$
components are complex numbers lies in the space $\mathbb{C}^n$

\begin{Definition}{Subspace}{subspace}
    A \textbf{subspace}\footnote{\href{https://trello.com/c/qHJeDNkU}{Introduction to Linear Algebra, Strang, 4th Edition, 2009}, p. 122}
    of a \hyperlink{vector-space}{vector space} is a set of vectors (including $\boldsymbol{0}$) that satisfies 2
    requirements: If $\boldsymbol{v}$ and $\boldsymbol{w}$ are vectors in the subspace and $c$ any scalar, then

    \begin{itemize}
        \item $\boldsymbol{v} + \boldsymbol{w}$ is in the subspace
        \item $c\boldsymbol{v}$ is in the subspace
    \end{itemize}

    In other words, the set of vectors is ``\textbf{closed}" under addition and multiplication - \textbf{all linear combinations
    stay in the subspace}

\end{Definition}

Intuitively, we can visualize a subspace in the 3-dimensional space $\mathbb{R}^3$. Choose a plane through the
origin $(0, 0, 0)$. That plane is a vector space in its own right. If we add two vectors in the plane, their sum is in
the plane; if we multiply an in-plane vector by $2$ or $-5$, it is still in the plane. This plane is a vector space
\textbf{inside $\mathbb{R}^3$} or is a subspace of the full vector space $\mathbb{R}^3$

\begin{Definition}{Column Space \phantomsection\hypertarget{column-space}}{column-space}
    The \textbf{column space}, $C(A)$, consists of all linear combinations of the columns, i.e. the combinations of all
    possible vectors $A\boldsymbol{x}$

    The subspece $C(A)$ is the ``\textbf{span}" of matrix $A$
\end{Definition}

\begin{Definition}{Basis for a Vector Space}{basis}
    A \textbf{basis}\footnote{\href{https://trello.com/c/qHJeDNkU}{Introduction to Linear Algebra, Strang, 4th Edition, 2009}, p. 172} for a vector space is a sequence of vectors with 2 properties

    \begin{itemize}
        \item The basis vectors are linearly independent, and
        \item they \hyperlink{column-space}{span} the space
    \end{itemize}
\end{Definition}

For a basis $\{\boldsymbol{v_1}, \ldots, \boldsymbol{v_d}\}$ of $V$

\begin{Definition}{Bilinear Form \phantomsection\hypertarget{bilinear-form}}{bilinear-form}
    Let $F$ be a field and $V$ a vector space over $F$. A \textit{bilinear form} on $V$ is a function
    $B: V \times V \rightarrow F$ that is linear in each variable when the other one is fixed. That is

    \begin{align}
        B(v + v', w) &= B(v, w) + B(v', w) \\
        B(cv, w)     &= cB(v, w)
    \end{align}

    for all $v, v', w \in V$ and $c \in F$, and

    \begin{align}
        B(v, w + w') &= B(v, w) + B(v, w') \\
        B(v, cw)     &= cB(v, w)
    \end{align}

    for all $v, v, w' \in V$ and $c \in F$

    We call $B$ \textit{symmetric} when

    \begin{equation}
        B(v, w) = B(w, v)
    \end{equation}

    for all $v, w \in V$\footnote{\href{https://kconrad.math.uconn.edu/blurbs/linmultialg/bilinearform.pdf}{Keith Conrad}}

    \begin{tcbraster}[
        raster columns=2,
        raster equal height
    ]
        \begin{Definition}{Linear Transformation}{linear-transformation}
            The transformation is \textbf{linear} if it meets these requirements for all $\boldsymbol{v}$ and $\boldsymbol{w}$\footnote{\href{https://trello.com/c/qHJeDNkU}{Introduction to Linear Algebra, Strang, 4th Edition, 2009}, p. 375}:

            \begin{equation}
                T(c\boldsymbol{v} + d\boldsymbol{w}) = cT(\boldsymbol{v}) + dT(\boldsymbol{w})
            \end{equation}

            for all $c$ and $d$
        \end{Definition}
        \begin{Definition}{Field}{field}
            A field is a set of elements in which a pair of operations called multiplication and addition is defined analogous
            to the operations of multiplication and addition in the real number system (which is itself an example of a
            field)\footnote{\href{https://projecteuclid.org/ebooks/notre-dame-mathematical-lectures/Chapter-I-Linear-Algebra/chapter/Chapter-I-Linear-Algebra/ndml/1175197044}{A. Fields, Chapter I: Linear Algebra}, Galois Theory: Lectures Delivered at the University of Notre Dame, Project Euclid}
        \end{Definition}
    \end{tcbraster}

    \begin{Definition}{Group}{group}
        A group\footnote{\href{https://www.maths.gla.ac.uk/~mwemyss/teaching/3alg1-7.pdf}{Introduction to Group Theory, Michael Wemyss}} is a non-empty set $G$ together with a rule that assigns to each pair $g$ and $h$ of elements of $G$
        and element $g * h$ such that
        \begin{itemize}
            \item[] \includegraphics[width=0.05\textwidth]{嘟嘟可.png}  $g * h \in G$, which we say $G$ is \textbf{closed} under $*$
            \item[] \includegraphics[width=0.05\textwidth]{嘟嘟可.png} $g * (h * k) = (g * h) * k$ for all $g, h, k \in G$, which we call $*$ being
            \textbf{associative}
            \item[] \includegraphics[width=0.05\textwidth]{嘟嘟可.png} There exists an \textbf{identity} element $e \in G$ such that $e * g = g * e$ for all
            $g \in G$
            \item[] \includegraphics[width=0.05\textwidth]{嘟嘟可.png} Every element $g \in G$ has an \textbf{inverse} $g^{-1}$ such that
            $g * g^{-1} = g^{-1} * g = e$
        \end{itemize}
    \end{Definition}

    \begin{Definition}{Vector Space (Field Theory)}{vector-space-field}
        If $V$ is an additive abelian group with elements $A, B, \ldots$, $F$ a field with elements $a, b, \dots$, and if
        for each $a \in F$ and $A \in V$ the product $aA$ denotes element of $V$, then $V$ is called a
        \textbf{(left) vector space over F} if the following assumptions hold:

        \begin{align}
            a(A + B) &= aA + aB \\
            (a + b)A &= aA + bA \\
            a(bA) &= (ab)A \\
            1A &= A \\
        \end{align}
    \end{Definition}
\end{Definition}